\clearpage
\section{Normalizing the unchanged constant-result program}
\label{sec:proposition-result-normalization}
The first milestone in this section completes the rightmost structural link
of Figure~22. Let $P_R$ denote the \emph{original} inductive translation of
$\lambda f.\,(\lambda w.\,1)\,(f\,())$. Let $R_d$ denote the single large
right abstraction in the parametricity family, at allowance $d+2$.
The theorem proves an exact-count contextual equivalence between the padded
closed $P_R$ and $R_d$, for every $d\geq0$. It then connects the original
translation in arbitrary finite variable and atom environments to $R_d$
with genuine outside weakenings. No term, translation clause, primitive
rule or observation class changes.

\paragraph{Correspondence with the paper.}
The \href{https://arxiv.org/html/1907.01257v5\#S12.SS3}{paper's Figure~22}
labels this final link by associativity, commutativity, idempotence and
box-weakening pulling. Lean first applies the already proved Idem/BPullW
weakening calculation to the unused variable under the inner lambda. It
then isolates the entire remaining finite sharing tree and applies the
proved general sharing-tree equivalence. It does not enumerate a separate
sequence of elementary associativity and commutativity moves for this tree.
That broader tree theorem already has proved input safety, forward and
converse robustness, and exact execution consequences. This is a difference
in proof organization, not an assumed graph equality or a simplified translation.

\subsection{Kernel-checked finite representation certificates}
\termdef{finite-incidence-certificate}{Finite incidence certificate}.
A finite incidence certificate consists of two candidate carrier maps and
their inverses, together with proofs preserving every internal wire label,
every recursive edge label, and every ordered source and target list. External
ports are fixed. An auxiliary table function looks up a carrier in an explicit
source inventory and returns the corresponding candidate target; its fallback
only makes the function total. Acceptance depends on the inverse and incidence
proofs, so the fallback cannot hide an omitted carrier.

The inventory lemma turns a predicate checked at every listed entry into a
predicate on the whole carrier, using the inventory's completeness proof.
An empty inventory admits no element; applying this observation twice gives
the unique bijection between two empty carriers. For nonempty carriers, the
bijection constructor uses the two checked inverse laws. The isomorphism
constructor applies inventory completeness to all label and ordered-incidence
checks. These arguments neither assume that inventories have no repeats nor
infer graph equality from matching counts.
\lean{Copying/FiniteGraphCertificate}{23}{24}{FiniteGraphCertificate.table}
\lean{Copying/FiniteGraphCertificate}{26}{27}{FiniteGraphCertificate.all}
\lean{Copying/FiniteGraphCertificate}{29}{36}{FiniteGraphCertificate.emptyBijection}
\lean{Copying/FiniteGraphCertificate}{37}{44}{FiniteGraphCertificate.bijection}
\lean{Copying/FiniteGraphCertificate}{45}{59}{FiniteGraphCertificate.iso}

The edge-rank lemma checks that every vertex shared by the target of one edge
and the source of another makes the rank strictly increase. Completeness of
the two edge inventories supplies this inequality for every adjacency.
The existing path-rank theorem then shows that every finite path has distinct
edges. Thus a finite computation establishes acyclicity without assuming it.
\lean{Copying/FiniteGraphCertificate}{60}{65}{FiniteGraphCertificate.path_nodup}

The small carrier-equality tactic recursively handles the actual finite sums
and dependent tagged components. All concrete checks use Lean's kernel-checked
\texttt{decide}; no native-evaluation axiom, external graph matcher, or
unverified isomorphism oracle enters the proof. Candidate tables are proof
inputs, and every certificate is rechecked against the original graph.
\lean{Copying/FiniteGraphCertificate}{1}{22}{Finite carrier equality and certificate setup}

\subsection{The literal graph and its complete sharing component}
The two scoped terms name the callback $f\,()$ and the outer lambda's body
$(\lambda w.\,1)\,(f\,())$, respectively, in the one-variable environment.
They are subterms of the unchanged right program. The mixed intermediate
keeps the original application and callback translation but replaces the
inner lambda translation by its normalized closed abstraction beside an
outside term weakening. Its callback is padded to allowance one exactly
as prescribed by the original translation.
\lean{Copying/PropositionResultGraphs}{11}{11}{Result.callbackTerm}
\lean{Copying/PropositionResultGraphs}{12}{12}{Result.bodyTerm}
\lean{Copying/PropositionResultGraphs}{14}{19}{Result.mixed}

\termdef{result-sharing-tree}{Result sharing tree}.
Write $W$ for a term weakening and $C_m$ for the original canonical
$m$-input term tree. The complete retained sharing component is
\[
 T=\bigl(W\otimes((C_1\otimes W);C_2)\bigr);C_2.
\]
It contains five weakenings and five contractions. The first inner $W$
comes from the unused variable of the unit argument, and the outer $W$
comes from the unused variable of the closed sequencing lambda. Both
original two-use distributors retain their own weakening. No canonical
row is silently replaced by a binary contraction.

For any one-input, one-output graph $U$, the schematic body is the outer
application of the closed constant-returning lambda to the callback
application, with $U$ on the callback's function wire and the literal unit
on its argument wire. Substituting $T$ gives the mixed graph's actual
sharing structure; substituting $C_1$ gives the original schematic
\texttt{callOneBody}.
\lean{Copying/PropositionResultGraphs}{20}{24}{Result.sharing}
\lean{Copying/PropositionResultGraphs}{26}{30}{Result.schematic}

The first finite isomorphism identifies the original constant translation
in one variable with the schematic numeral beside an ignored binder. Both
have exactly the same two edges; the certificate exchanges their carrier
order and fixes the term input and unused output. The second isomorphism
factors the whole mixed graph into the schematic operational skeleton and
$T$. Its fifteen-edge certificate retains both application heads, the unit,
the lambda, its entire original box body and all ten sharing edges. It
checks every ordered incidence, including both contraction inputs.
\lean{Copying/PropositionResultGeometry}{13}{35}{Result.ConstantBody.iso}
\lean{Copying/PropositionResultGeometry}{39}{65}{Result.SharingFactorization.iso}

\subsection{Pulling the unused resource through the real lambda}
The closed inner lambda is first identified with the small schematic right
abstraction: remove only the empty outer distributor and apply the numeral
body isomorphism under its original box. The full term-weakening theorem
then connects the actual inner lambda in environment $f$ to that closed
abstraction with one outside weakening. Its only semantic leaves are the
already established Idem and BPullW certificates.
\lean{Copying/PropositionResultNormalization}{14}{19}{Result.closedSmall_iso}
\lean{Copying/PropositionResultNormalization}{20}{26}{Result.inner_weakening}

Lift this calculation through the left operand of the unchanged outer
application. The padded callback operand and both application distributors
remain present. Lift the result through the original outer abstraction,
and remove its genuinely empty outside distributor by an isomorphism.
The resulting calculation runs from the unchanged closed $P_R$ to the
lambda payload containing the mixed body.
\lean{Copying/PropositionResultNormalization}{27}{31}{Result.body_weakening}
\lean{Copying/PropositionResultNormalization}{32}{42}{Result.program_weakening}

\subsection{Replacing the whole sharing tree under both applications}
The proof of the sharing-tree predicate exhausts all ten edges, showing that
each is a term contraction or term weakening. In the graph's original edge
enumeration the rank list is
\[
 (0,0,1,0,0,3,2,0,5,4).
\]
Every actual adjacency increases this rank, as checked against the source
and target lists. The finite rank lemma therefore establishes acyclicity.
The resulting actual sharing-tree predicate, together with the existing
canonical-tree predicate, supplies a genuine sharing-template relation
between $T$ and $C_1$ and hence a recursive filling witness.
\lean{Copying/PropositionResultNormalization}{43}{56}{Result.sharing_tree}
\lean{Copying/PropositionResultNormalization}{57}{62}{Result.sharing_witness}

The body witness retains the unit, inner closed lambda and two application
heads by reflexive fillings, tensors and serially composes them with the
sharing witness, and changes the left representation using the whole-graph
factorization certificate. Its right endpoint is definitionally the
schematic call-one body. Boxing this witness with one bound variable and
no free resources, and retaining the outer lambda edge, produces the
program witness. The selected bound output is preserved throughout.
\lean{Copying/PropositionResultNormalization}{63}{76}{Result.body_witness}
\lean{Copying/PropositionResultNormalization}{77}{81}{Result.program_witness}

The original, weakened and normalized open-program records package precisely
these three graphs with proved focus and hole freedom. The original record
uses translation invariants. The weakened record obtains both invariants
from the actual weakening calculation. The normalized record uses the
schematic family's recursive label theorem. None of these records contains
a field assuming observational equivalence.
\lean{Copying/PropositionResultNormalization}{83}{84}{Result.original}
\lean{Copying/PropositionResultNormalization}{85}{86}{Result.weakened}
\lean{Copying/PropositionResultNormalization}{87}{93}{Result.normalized}

The first contextual comparison is the exact-count soundness of the
weakening calculation. For the second, fill the sharing witness into each
original binding-free closing context and apply the completed sharing
robustness theorem. Both closed endpoints are hole-free by their actual
program and context certificates. Composition preserves exact successful
transition counts in both directions. The termination equivalence is
obtained only afterwards by forgetting equality of costs.
\lean{Copying/PropositionResultNormalization}{94}{100}{Result.original_normalized_exact}
\lean{Copying/PropositionResultNormalization}{102}{104}{Result.original_normalized}

\subsection{All adequate depth allowances}
Padding is applied to structural certificates, rather than inferred from
determinism. For any padding-closed template, the witness-padding helper
writes $N=d+n$ and uses the proved $n$-fold padding construction, transporting
the endpoints through the original padding equalities. The large-right
padding isomorphism follows by induction on the extra allowance: the base
case is self-padding, and the successor uses padding congruence followed
by the previously proved one-layer schematic-abstraction isomorphism.
\lean{Copying/PropositionResultPadding}{11}{20}{Result.padWitness}
\lean{Copying/PropositionResultPadding}{21}{28}{Result.largeRight_padding}

The three allowance-indexed program records pad the original and weakened
graphs to $d+2$, while the normalized record uses the literal schematic
abstraction at that allowance. The padded program witness keeps its left
endpoint and uses the large-right padding isomorphism on its right endpoint.
\lean{Copying/PropositionResultPadding}{30}{30}{Result.originalAt}
\lean{Copying/PropositionResultPadding}{31}{31}{Result.weakenedAt}
\lean{Copying/PropositionResultPadding}{32}{34}{Result.normalizedAt}
\lean{Copying/PropositionResultPadding}{36}{43}{Result.program_witnessAt}

The same two-stage soundness proof now applies at every allowance $d+2$:
pad the original weakening calculation, then apply the padded sharing
witness in an arbitrary closing context. Exact costs remain equal. Only
the separate termination corollary forgets them. In particular, allowance
three, needed by the allocating left-hand program, is already covered.
\lean{Copying/PropositionResultPadding}{44}{52}{Result.original_normalizedAt_exact}
\lean{Copying/PropositionResultPadding}{54}{57}{Result.original_normalizedAt}

\subsection{Arbitrary variable and atom environments}
Padding commutes with actual outside weakening. Induct first over the
variable block and then over the atom block, applying the established
single-extension padding equation at each successor. This proves equality
of the two assembled graphs, including their original boundary order.
\lean{Copying/PropositionResultEnvironments}{10}{22}{Result.ambient_padTo}

The original-environment program is exactly \texttt{translatedProgram} of
the unchanged constant-result term at allowance $d+2$. The normalized
program is the schematic abstraction with its certified ambient weakenings.
Pad the full term-weakening calculation for the original environment, then
use the commuting equation to identify its endpoint with the actual ambient
extension of the padded closed translation. This first link is exact-count.
\lean{Copying/PropositionResultEnvironments}{24}{25}{Result.originalEnvironment}
\lean{Copying/PropositionResultEnvironments}{26}{30}{Result.normalizedEnvironment}
\lean{Copying/PropositionResultEnvironments}{31}{39}{Result.environment_weakening_exact}

Compose this environment link with the closed normalization lifted through
the already proved sole-hole ambient context. The resulting termination
comparison holds for every finite pair of environment sizes. The named
version substitutes the lengths of the original distinct lists. The
observation corollary quantifies over every original binding-free closing
context, and states equivalence of successful termination; it assumes
neither callback termination nor single use of the returned closure.
\lean{Copying/PropositionResultEnvironments}{40}{47}{Result.original_normalized_environment}
\lean{Copying/PropositionResultEnvironments}{49}{56}{Result.original_normalized_named}
\lean{Copying/PropositionResultEnvironments}{57}{60}{Result.original_normalized_termination}
